\documentclass[10pt,a4paper]{amsart}
\usepackage{amsmath,amssymb,amsthm,graphicx}
\usepackage{verbatim}
\newenvironment{reference}{\comment}{\endcomment}
%\newenvironment{reference}{}{}
\newenvironment{slogan}{\comment}{\endcomment}
\newenvironment{history}{\comment}{\endcomment}

% For commutative diagrams we use Xy-pic
\usepackage[all]{xy}

% We use 2cell for 2-commutative diagrams.
\xyoption{2cell}
% Theorem environments.
%
\theoremstyle{plain}
\newtheorem{theorem}[subsection]{Theorem}
\newtheorem{proposition}[subsection]{Proposition}
\newtheorem{lemma}[subsection]{Lemma}

\theoremstyle{definition}
\newtheorem{definition}[subsection]{Definition}
\newtheorem{example}[subsection]{Example}
\newtheorem{exercise}[subsection]{Exercise}
\newtheorem{situation}[subsection]{Situation}

\theoremstyle{remark}
\newtheorem{remark}[subsection]{Remark}
\newtheorem{remarks}[subsection]{Remarks}

\numberwithin{equation}{subsection}

% Macros
%
\def\lim{\mathop{\mathrm{lim}}\nolimits}
\def\colim{\mathop{\mathrm{colim}}\nolimits}
\def\Spec{\mathop{\mathrm{Spec}}}
\def\Hom{\mathop{\mathrm{Hom}}\nolimits}
\def\Ext{\mathop{\mathrm{Ext}}\nolimits}
\def\SheafHom{\mathop{\mathcal{H}\!\mathit{om}}\nolimits}
\def\SheafExt{\mathop{\mathcal{E}\!\mathit{xt}}\nolimits}
\def\Sch{\mathit{Sch}}
\def\Mor{\mathop{\mathrm{Mor}}\nolimits}
\def\Ob{\mathop{\mathrm{Ob}}\nolimits}
\def\Sh{\mathop{\mathit{Sh}}\nolimits}
\def\NL{\mathop{N\!L}\nolimits}
\def\CH{\mathop{\mathrm{CH}}\nolimits}
\def\proetale{{pro\text{-}\acute{e}tale}}
\def\etale{{\acute{e}tale}}
\def\QCoh{\mathit{QCoh}}
\def\Ker{\mathop{\mathrm{Ker}}}
\def\Im{\mathop{\mathrm{Im}}}
\def\Coker{\mathop{\mathrm{Coker}}}
\def\Coim{\mathop{\mathrm{Coim}}}

% Boxtimes
%
\DeclareMathSymbol{\boxtimes}{\mathbin}{AMSa}{"02}

%
% Macros for moduli stacks/spaces
%
\def\QCohstack{\mathcal{QC}\!\mathit{oh}}
\def\Cohstack{\mathcal{C}\!\mathit{oh}}
\def\Spacesstack{\mathcal{S}\!\mathit{paces}}
\def\Quotfunctor{\mathrm{Quot}}
\def\Hilbfunctor{\mathrm{Hilb}}
\def\Curvesstack{\mathcal{C}\!\mathit{urves}}
\def\Polarizedstack{\mathcal{P}\!\mathit{olarized}}
\def\Complexesstack{\mathcal{C}\!\mathit{omplexes}}
% \Pic is the operator that assigns to X its picard group, usage \Pic(X)
% \Picardstack_{X/B} denotes the Picard stack of X over B
% \Picardfunctor_{X/B} denotes the Picard functor of X over B
\def\Pic{\mathop{\mathrm{Pic}}\nolimits}
\def\Picardstack{\mathcal{P}\!\mathit{ic}}
\def\Picardfunctor{\mathrm{Pic}}
\def\Deformationcategory{\mathcal{D}\!\mathit{ef}}
\usepackage{cleveref}
\crefname{originalreference}{original reference}{original references}
\Crefname{originalreference}{Original reference}{Original references}
\makeatletter
\newlabel{lemma-strange-regular-sequence}{{0640}{source}{}{sourcepointer-1}{}}
\newlabel{lemma-strange-regular-sequence@cref}{{[originalreference][0][]0640}{[1][0][]source}{}{}{}}
\makeatother
\crefname{originallemma}{original lemma}{original lemmas}
\Crefname{originallemma}{Original lemma}{Original lemmas}
\makeatletter
\newlabel{more-algebra-lemma-regular-koszul-regular}{{062F}{source}{}{sourcepointer-2}{}}
\newlabel{more-algebra-lemma-regular-koszul-regular@cref}{{[originallemma][0][]062F}{[1][0][]source}{}{}{}}
\makeatother
\title{Flat base change can destroy the Koszul-regularity of a regular sequence}
\author{The Stacks Project Authors}
\date{stacks-project:examples.tex:local-snapshot-20261012; GNU FDL 1.2 or later}
\begin{document}
\maketitle
\noindent\textit{Editorial note.} All delivered mathematical text is the author's own, copied from the pinned Stacks Project chapter file examples.tex; no statement or proof step was rewritten, reordered or removed.\par\medskip
\noindent\textit{Editorial note.} The counted proof is the author's complete argument: Companion of lemma-strange-regular-sequence, which supplies the ring B and its regular sequence.\par\medskip
\noindent\textit{Editorial note.} The label lemma-base-change-regular-sequence names the counted statement, and the proof environment immediately following it in the source is delivered verbatim as its proof.\par\medskip
\noindent\textit{Editorial note.} The counted statement is carried by the author's own lemma environment in the chapter source.\par\medskip
\noindent\textit{Allowed standard prerequisites (editorial).}\begin{itemize}
\item The module E and the square-zero extension ring of the section construction, whose properties (a) through (d) are proved in section prose and restated in the companion lemma lemma-strange-regular-sequence
\item Over a principal ideal domain a finitely generated torsion-free module is flat; the ring k[z] localized at the prime generated by z is a discrete valuation ring
\item More on Algebra, Lemma more-algebra-lemma-regular-koszul-regular: a regular sequence is Koszul-regular; hence a pair whose Koszul homology does not vanish cannot be a regular sequence
\end{itemize}\medskip
\ifdefined\hypertarget\hypertarget{sourcepointer-1}{}\fi
\noindent\textit{Source reference.} Original reference 0640: sources/stacks-examples/pinned-original-source.tex, source lines 1333--1333. Prerequisite pointer invoked at 1362; it names a standard fact of the same chapter that the counted argument uses, and every step of the counted proof built on it is retained verbatim in the delivered body.\par\medskip
\ifdefined\hypertarget\hypertarget{sourcepointer-2}{}\fi
\noindent\textit{Source reference.} Original lemma 062F: sources/stacks-chow/tags.txt, source lines 7794--7794. Prerequisite pointer invoked at 1371; it names a standard fact of a sibling chapter that the counted argument uses, and every step of the counted proof built on it is retained verbatim in the delivered body.\par\medskip
\begin{lemma}
\label{lemma-base-change-regular-sequence}
There exists a local homomorphism of local rings $A \to B$
and a regular sequence $x, y$ in the maximal ideal of $B$ such that
$B/(x, y)$ is flat over $A$, but such that the images
$\overline{x}, \overline{y}$ of $x, y$ in $B/\mathfrak m_AB$ do not
form a regular sequence, nor even a Koszul-regular sequence.
\end{lemma}

\begin{proof}
Set $A = k[z]_{(z)}$ and let $B = (k[x, y, z] \oplus E)_{(x, y, z, E)}$.
Since $x, y, z$ is a regular sequence in $B$, see proof of
Lemma \ref{lemma-strange-regular-sequence},
we see that $x, y$ is a regular sequence in $B$ and that
$B/(x, y)$ is a torsion free $A$-module, hence flat.
On the other hand, there exists a nonzero element
$\delta \in B/\mathfrak m_AB = B/zB$ which is annihilated
by $\overline{x}, \overline{y}$. Hence
$H_2(K_\bullet(B/\mathfrak m_AB, \overline{x}, \overline{y})) \not = 0$.
Thus $\overline{x}, \overline{y}$ is not Koszul-regular, in particular
it is not a regular sequence, see
More on Algebra, Lemma \ref{more-algebra-lemma-regular-koszul-regular}.
\end{proof}

\end{document}
